% Figure: convolution of a rectangular input pulse with the RC
% impulse response, showing the input, impulse response, and the
% output (a charge-discharge transient).
\begin{figure}[htbp]
\centering
\begin{tikzpicture}

% Top panel: input pulse v_in(t) = 1 on [0, 1], 0 elsewhere
\begin{axis}[
  name=in,
  width=11cm, height=3.5cm,
  xlabel={},
  ylabel={$v_{\mathrm{in}}(t)$},
  xmin=-0.3, xmax=3.5,
  ymin=-0.2, ymax=1.3,
  axis lines=left,
  grid=both,
  major grid style={engblack!15},
  font=\small,
  xticklabels={},
  title={Input pulse},
]
\addplot[engblue, very thick]
  coordinates {(-0.3, 0) (0, 0) (0, 1) (1, 1) (1, 0) (3.5, 0)};
\end{axis}

% Middle panel: impulse response h(t) = (1/RC) exp(-t/RC), with RC = 0.5
\begin{axis}[
  at={(in.south)}, anchor=north, yshift=-8pt,
  name=imp,
  width=11cm, height=3.5cm,
  xlabel={},
  ylabel={$h(t)$},
  xmin=-0.3, xmax=3.5,
  ymin=-0.2, ymax=2.2,
  domain=0:3.5,
  samples=200,
  axis lines=left,
  grid=both,
  major grid style={engblack!15},
  font=\small,
  xticklabels={},
  title={Impulse response (RC, $\tau = 0.5$)},
]
\addplot[draw=engblack, very thick] {1/0.5 * exp(-x/0.5)};
\end{axis}

% Bottom panel: output. For a pulse of unit height on [0, 1] passing
% through an RC filter with tau = 0.5, the output is:
%   v_out(t) = 1 - exp(-t/tau)               for 0 <= t <= 1
%   v_out(t) = (1 - exp(-1/tau)) exp(-(t-1)/tau)  for t > 1
\begin{axis}[
  at={(imp.south)}, anchor=north, yshift=-8pt,
  width=11cm, height=4cm,
  xlabel={$t$},
  ylabel={$v_{\mathrm{out}}(t)$},
  xmin=-0.3, xmax=3.5,
  ymin=-0.2, ymax=1.3,
  axis lines=left,
  grid=both,
  major grid style={engblack!15},
  font=\small,
  title={Output: $v_{\mathrm{out}} = v_{\mathrm{in}} * h$},
]
% Charging phase
\addplot[engred, very thick, domain=0:1, samples=100]
  {1 - exp(-x/0.5)};
% Discharging phase
\addplot[engred, very thick, domain=1:3.5, samples=100]
  {(1 - exp(-1/0.5)) * exp(-(x-1)/0.5)};
% Faint extension of the pulse for visual alignment
\addplot[engblue!30, very thick]
  coordinates {(-0.3, 0) (0, 0) (0, 1) (1, 1) (1, 0) (3.5, 0)};
\end{axis}

\end{tikzpicture}
\caption[Convolution of a pulse with the RC impulse
response]{Convolution of a unit-height rectangular input pulse on
$[0, 1]$ (top) with the RC impulse response $h(t) = (1/\tau)
e^{-t/\tau}$ for $\tau = 0.5$ (middle) gives the charge-discharge
output (bottom). On $[0, 1]$ the output rises as $1 - e^{-t/\tau}$
while the input holds the capacitor at unit drive; at $t = 1$ the
input drops to zero and the output decays from its peak as
$(1 - e^{-1/\tau}) e^{-(t-1)/\tau}$. The convolution integral
\textcite{text:oppenheim-schafer} reproduces the time-domain
solution that a Volume III circuit analysis derives from the
governing ODE.}
\label{fig:vol02:ch06:rc-convolution}
\end{figure}
